\chapter{Charte d'utilisation de l'IA générative}

\section*{1. Usages dans le travail scientifique}

L'ensemble des choix scientifiques, des raisonnements, des développements théoriques et des implémentations relève de mon travail personnel. Les outils d'IA générative n'ont pas été utilisés pour produire directement~:

\begin{itemize}
  \item les questions de recherche~;
  \item les hypothèses~;
  \item la méthodologie expérimentale~;
  \item la conception des expériences~;
  \item l'analyse ou l'interprétation des résultats~;
  \item la modélisation ou la formalisation théorique~;
  \item les implémentations soumises~;
  \item les contributions scientifiques originales présentées dans ce travail.
\end{itemize}

\section*{2. Usages dans la rédaction du manuscrit}

Les outils d'IA générative n'ont pas été utilisés pour générer substantiellement le contenu scientifique, l'argumentation, l'état de l'art, l'analyse des résultats ou les conclusions.

\section*{3. Usage(s) spécifique(s) de l'IA générative autorisé(s) par le promoteur dans le cadre du mémoire}

%\textit{(à compléter par le promoteur)}

En accord avec le promoteur, l'IA générative a été utilisée comme 
outil d'assistance pour les tâches suivantes :

\begin{itemize}
    \item \textbf{Génération de visuels} : production de spécifications PlantUML pour les diagrammes d'architecture (figures du chapitre Architecture), explicitement référencées par une note méthodologique dédiée. Tous les diagrammes ont été manuellement relus et corrigés pour assurer leur fidélité au code source de MetaDONE.
    \item \textbf{Assistance rédactionnelle} : suggestions de reformulation pour améliorer la clarté, la concision et la cohérence stylistique du texte français. Les choix de structure, d'argumentation et de message scientifique relèvent du travail personnel de l'auteure ; l'assistance n'a porté que sur la forme.
    \item \textbf{Assistance au codage} : aide ponctuelle à la rédaction du code de benchmark (scénarios 1, 2, 3, 5) et à la mise en forme des correctifs d'optimisation. L'analyse des goulots d'étranglement, le diagnostic des causes racines et la conception des optimisations relèvent du travail personnel ; je comprends et maîtrise chaque ligne du code commité.
\end{itemize}


\section*{4. Responsabilité}

Je confirme être l'auteur des contributions scientifiques présentées dans ce mémoire et être capable d'expliquer et de défendre l'ensemble de son contenu. Aucune utilisation d'IA générative n'a remplacé mon travail intellectuel.

\vspace{2cm}

\noindent\begin{tabular}{p{6cm}p{6cm}}
\textit{Signature de l'étudiante} & \textit{Signature du promoteur} \\[3cm]
\hrulefill & \hrulefill \\
Carine Pochet & Prof. Vincent Englebert \\
\end{tabular}
